\section{F8 · Fouriertransformen}

\textbf{Kjernen:} Fourierrekker for funksjoner som ikke er periodiske, ved å la $L\to\infty$. Summen over $k$ blir et integral over frekvensen $\omega$. Derivasjon blir multiplikasjon med $\ii\omega$, og konvolusjon blir produkt. Det gjør PDE-er på hele $\R$ om til ODE-er (F10).

\begin{formel}[Fouriertransform og invers\fs]
For $f\in L^1(\R)$ (det vil si $\int_{-\infty}^\infty\abs f<\infty$):
\begin{gather*}
\hat f(\omega)=\Four(f)(\omega)=\frac1{\sqrt{2\pi}}\int_{-\infty}^{\infty}f(x)\ee^{-\ii\omega x}\dd x,\\
f(x)=\Four^{-1}(\hat f)(x)=\frac1{\sqrt{2\pi}}\int_{-\infty}^{\infty}\hat f(\omega)\ee^{\ii\omega x}\dd\omega .
\end{gather*}
Inversjonen gjelder når både $f$ og $\hat f$ er absolutt integrerbare. $f$ kalles tidsdomenet og $\hat f$ frekvensdomenet.
\end{formel}
\begin{fallgruve}[konvensjoner]
Andre bøker plasserer $2\pi$ annerledes: $\int f\ee^{-\ii\omega x}$ (med $\frac1{2\pi}$ i inversen), $\frac1{2\pi}\int\dots$ eller $\int f\ee^{-2\pi\ii\omega x}$ (frekvens i Hz), og noen har minuset i inversen. Sjekk alltid hvilken definisjon en tabell eller en gammel eksamen bruker. Kurset bruker $\frac1{\sqrt{2\pi}}$ begge veier.
\end{fallgruve}

\begin{defn}[cosinus- og sinustransform]
For like $f$: $\hat f_c(\omega)=\sqrt{\frac2\pi}\int_0^\infty f(x)\cos(\omega x)\dd x$, med invers $\check g_c(x)=\sqrt{\frac2\pi}\int_0^\infty g(\omega)\cos(\omega x)\dd\omega$.\\
For odde $f$: $\hat f_s(\omega)=\sqrt{\frac2\pi}\int_0^\infty f(x)\sin(\omega x)\dd x$, med invers på samme form.\\
En reell funksjon gir en reell transform. Sammenhengen er $\hat f=\hat f_c$ for like $f$ og $\hat f=-\ii\hat f_s$ for odde $f$.
\end{defn}

\begin{formel}[Transformtabellen\fs]
\begin{center}
\begin{tabular}{@{}ll@{}}
\toprule
$f(x)$ & $\hat f(\omega)$\\
\midrule
$\ee^{-ax^2}$ \ ($a>0$) & $\dfrac1{\sqrt{2a}}\ee^{-\omega^2/(4a)}$ \quad (gaussisk gir gaussisk)\\[6pt]
$\ee^{-a\abs x}$ \ ($a>0$) & $\sqrt{\dfrac2\pi}\,\dfrac{a}{\omega^2+a^2}$\\[6pt]
$1$ for $\abs x<a$, ellers $0$ & $\sqrt{\dfrac2\pi}\,\dfrac{\sin(a\omega)}{\omega}$\\[6pt]
$\sqrt{\dfrac2\pi}\,\dfrac{\sin(ax)}x$ & $1$ for $\abs\omega<a$, ellers $0$\\
\bottomrule
\end{tabular}
\end{center}
Ikke på arket, men nyttig: $\chi_{[a,b]}$ (1 på $[a,b]$, ellers 0) har $\hat f(\omega)=\dfrac{\ee^{-\ii\omega a}-\ee^{-\ii\omega b}}{\ii\omega\sqrt{2\pi}}$ for $\omega\neq0$ og $\hat f(0)=\dfrac{b-a}{\sqrt{2\pi}}$.
\end{formel}

\begin{teorem}[regneregler (de tre første er fra forelesningen)]
\begin{itemize}
  \item \textbf{Linearitet:} $\Four(af+bg)=a\Four(f)+b\Four(g)$.
  \item \textbf{Derivasjon:} hvis $f,f'\in L^1$ og $f\to0$ for $x\to\pm\infty$, er $\Four(f')=\ii\omega\,\Four(f)$, og dermed $\Four(f'')=-\omega^2\Four(f)$.
  \item \textbf{Konvolusjon:} $(f*g)(x)=\displaystyle\int_{-\infty}^{\infty}f(y)g(x-y)\dd y$, med $f*g=g*f$ og
  \[
  \Four(f*g)=\sqrt{2\pi}\;\Four(f)\,\Four(g).
  \]
  Konvolusjon i tid blir multiplikasjon i frekvens.
\end{itemize}
\textbf{Ekstra (ikke på lysarkene, kontrollregnet):} forskyvning $\Four(f(x-a))=\ee^{-\ii\omega a}\hat f(\omega)$, modulasjon $\Four(\ee^{\ii ax}f)=\hat f(\omega-a)$, skalering $\Four(f(bx))=\frac1{\abs b}\hat f(\omega/b)$ og $\Four(xf(x))=\ii\,\frac{\dd}{\dd\omega}\hat f(\omega)$.
\end{teorem}

\begin{teorem}[gaussiske integraler (øving 8, oppg.~5)]
$\displaystyle\int_{-\infty}^\infty\ee^{-x^2}\dd x=\sqrt\pi$ (kvadrer og bytt til polarkoordinater), og
$\displaystyle\int_{-\infty}^\infty\ee^{-ax^2+bx}\dd x=\sqrt{\frac\pi a}\,\ee^{b^2/(4a)}$ for $a>0$. $\ee^{-x^2/2}$ er sin egen transform.
\end{teorem}

\begin{oppskrift}[typiske oppgaver]
\begin{enumerate}
  \item \textbf{FT fra definisjonen:} for stykkevise funksjoner, integrer over intervallene der $f\neq0$ og bruk delvis integrasjon for $x^n$-ledd. Behandle $\omega=0$ for seg. Sjekk med symmetri: like $f$ gir reell og like $\hat f$, og odde $f$ gir rent imaginær og odde $\hat f$.
  \item \textbf{Med tabell og regler:} skriv $f$ som en kombinasjon av tabellfunksjoner. For eksempel gir $t^2\ee^{-t^2}$ at du deriverer $\ee^{-t^2}$ to ganger, eller bruker $\Four(x^2f)=-\hat f''$.
  \item \textbf{Konvolusjon:} $f*g=\Four^{-1}\bigl(\sqrt{2\pi}\hat f\hat g\bigr)$. For eksempel:
  \begin{itemize}
    \item $\ee^{-x^2}*\ee^{-2x^2}=\sqrt{\pi/3}\,\ee^{-2x^2/3}$ (øving 7, oppg.~5, kontrollregnet).
    \item $\operatorname{sinc}*\operatorname{sinc}=\pi\operatorname{sinc}$ med $\operatorname{sinc}x=\frac{\sin x}x$ (øving 8, oppg.~4).
    \item $\chi_{(0,2)}*\chi_{(0,2)}$ er en trekantfunksjon (øving 8, oppg.~2).
  \end{itemize}
\end{enumerate}
\end{oppskrift}

\begin{oving}
\begin{itemize}[leftmargin=5.6em,labelwidth=5em,labelsep=0.6em,align=left]
\item[\lett] \href{https://github.com/gustavbv/Matte4C---eksamen/blob/main/oppgavesett/F08-fouriertransform/2-Bergmann-eksamener/010-TMA4125-30-35-2022-kont-oppg7-oppgave.pdf}{Matte 4N kont.\ 2022, oppg.~7} {\color{graa}(\href{https://github.com/gustavbv/Matte4C---eksamen/blob/main/oppgavesett/F08-fouriertransform/2-Bergmann-eksamener/010-TMA4125-30-35-2022-kont-oppg7-svar.pdf}{LF})}: konvolusjon av gaussfunksjoner med tabellen
\item[\lett] \href{https://github.com/gustavbv/Matte4C---eksamen/blob/main/oppgavesett/F08-fouriertransform/3-hele-oppgaven-i-pensum/046-TMA4135-2015-host-oppg2-oppgave.pdf}{TMA4135 høst 2015, oppg.~2} {\color{graa}(\href{https://github.com/gustavbv/Matte4C---eksamen/blob/main/oppgavesett/F08-fouriertransform/3-hele-oppgaven-i-pensum/046-TMA4135-2015-host-oppg2-svar.pdf}{LF})}: $f_a*f_b$ for gaussfunksjoner
\item[\middels] \href{https://github.com/gustavbv/Matte4C---eksamen/blob/main/oppgavesett/F08-fouriertransform/2-Bergmann-eksamener/009-TMA4125-2022-var-oppg6-oppgave.pdf}{Matte 4N vår 2022, oppg.~6} {\color{graa}(\href{https://github.com/gustavbv/Matte4C---eksamen/blob/main/oppgavesett/F08-fouriertransform/2-Bergmann-eksamener/009-TMA4125-2022-var-oppg6-svar.pdf}{LF})}: FT av en avkortet cosinus
\item[\middels] \href{https://github.com/gustavbv/Matte4C---eksamen/blob/main/oppgavesett/F08-fouriertransform/1-TMA4432-ovinger-og-proveeksamen/006-TMA4432-oving-8-oppg2-oppgave.pdf}{Øving 8, oppg.~2} {\color{graa}(\href{https://github.com/gustavbv/Matte4C---eksamen/blob/main/oppgavesett/F08-fouriertransform/1-TMA4432-ovinger-og-proveeksamen/006-TMA4432-oving-8-oppg2-svar.pdf}{LF})}: $f*f$ fra definisjonen, og FT av $f$ og $g$
\item[\middels] \href{https://github.com/gustavbv/Matte4C---eksamen/blob/main/oppgavesett/F08-fouriertransform/3-hele-oppgaven-i-pensum/013-TMA4130-35-2025-kont-2-sett-oppg4-oppgave.pdf}{TMA4130 kont.\ 2025 (2.\ sett), oppg.~4} {\color{graa}(\href{https://github.com/gustavbv/Matte4C---eksamen/blob/main/oppgavesett/F08-fouriertransform/3-hele-oppgaven-i-pensum/013-TMA4130-35-2025-kont-2-sett-oppg4-svar.pdf}{LF})}: FT av stykkevis funksjon og derivasjonsregelen
\item[\vanskelig] \href{https://github.com/gustavbv/Matte4C---eksamen/blob/main/oppgavesett/F08-fouriertransform/1-TMA4432-ovinger-og-proveeksamen/005-TMA4432-oving-8-oppg1-oppgave.pdf}{Øving 8, oppg.~1} {\color{graa}(\href{https://github.com/gustavbv/Matte4C---eksamen/blob/main/oppgavesett/F08-fouriertransform/1-TMA4432-ovinger-og-proveeksamen/005-TMA4432-oving-8-oppg1-svar.pdf}{LF})}: FT av $x^2$ på $\abs x<1$ (mye delvis integrasjon)
\item[\vanskelig] \href{https://github.com/gustavbv/Matte4C---eksamen/blob/main/oppgavesett/F08-fouriertransform/3-hele-oppgaven-i-pensum/012-TMA4130-35-2025-kont-oppg7-oppgave.pdf}{TMA4130 kont.\ 2025, oppg.~7} {\color{graa}(\href{https://github.com/gustavbv/Matte4C---eksamen/blob/main/oppgavesett/F08-fouriertransform/3-hele-oppgaven-i-pensum/012-TMA4130-35-2025-kont-oppg7-svar.pdf}{LF})}: skaleringsregel og produkt som konvolusjon
\item[\vanskelig] \href{https://github.com/gustavbv/Matte4C---eksamen/blob/main/oppgavesett/F08-fouriertransform/3-hele-oppgaven-i-pensum/059-TMA4135-2011-kont-oppg2-oppgave.pdf}{TMA4135 kont.\ 2011, oppg.~2} {\color{graa}(\href{https://github.com/gustavbv/Matte4C---eksamen/blob/main/oppgavesett/F08-fouriertransform/3-hele-oppgaven-i-pensum/059-TMA4135-2011-kont-oppg2-svar.pdf}{LF})}: delbrøk, FT og et vanskelig integral
\end{itemize}
\end{oving}

